% Pursuing Stacks (cotes 134-2 to 134-8): archive cote 134-7, site page 107 (PDF page 108), typescript page 549.
% typescript. Draft compared: ps6.tex:5393-5452.
% AI-transcribed by Claude Opus 5.5 (claude-opus-5-5, adaptive thinking, effort high), Message Batches API, batch ps4 (msgbatch_012Aoc1NuskATPsyBeKSYK7J), 2026-11-12; not yet checked by eye.
\scanpage{134-7/107}
\mspage{549}

Our tentative $\overset{\hand{$\sim$}}{S}$ functor in (10) has been constructed in the most evi-%
dent way,  in order to satisfy an adjunction formula (11),  valid on
the level of semisimplicial objects  (and carr\xout{\unsure{a}}ying over\hand{,} hopefully\hand{,}
to the similar adjunction formula for \inserted[typed]{suitable}\xout{\illegible{2}} derived categories).
Next question is then,  what about the inversion formula
\[
(12)\qquad L\overset{\hand{$\sim$}}{H}.(\overset{\hand{$\sim$}}{S}(L_*))\simeq L_* \qquad ?
\]
T\over{\unsure{w}}{h}e question makes sense,  as  $L\overset{\hand{$\sim$}}{H}.$  is \xout{constructed (in} defined for any
semisimplicial affine scheme-with-section\inserted{over k}, or equivalently for any
co-semis\uline{i}mplicial augmented k-algebra $A^*$ over k,  by taking the augmentati%
on ideal  $\underline{m}^*$  in the latter and retaining only it's linear (co-semisimpli-%
cial) structure.  Keeping this in mind,  formula (12) comes out indeed as
a tautology again !

The \xout{dex} tentative description we just got is indeed of a most sedu-%
cing simplicity,  as seducing indeed as the description of homotopy in
terms of the Lie functor,  and closely related to the latter.  It gives
as a particular case an exceedingly simple description of the sought
four ``spheres over k''  $S(k,n)$.  But it is clear that this description is
liable to make sense only at the price of suitably extending the notion
of a ``bundle'' we are working with,  in a rather different direction I
would say from adding (or allowing) connected components, as suggested in
previous section.  Maybe we might view it though as a kindred,  but some-%
what subtler extension of our initial bundles,  namely that we are
now allowing,  not a discrete non-trivial ``set'' or k-scheme of connected
components,  but rather,  an infinitesimal one.  More specifically,
the suggestion which comes to mind here,  is to call now ``bundle'' \inserted{over k,} any
scheme  X over  k ,  admitting a subscheme
\[
X_o \subset X
\]
in such a way that $X_o$ should be a V-bundle (namely isomorphic to a vec-%
tor bundle) over k,  and X should be an infinitesimal neighbourhood of
$X_o$ ,  i.e. $X_o$ should be definable by a quasi-coherent Ideal  on X which
is nilpotent.  Equivalently,  in terms of the affine ring A of X,  we
are demanding that A should admit a nilpotent ideal J (which is of course
not part of it's structure),  such that A/J should be isomorphic to a
symmetric algebra over k (with respect to some k-module M).  Possibly
we may have moreover to impose further flatness restrictions.

When working with this extended notion of bundles,  there is no
problem for describing,  for the corresponding semisimplicial models

\margin{107}
